\documentclass{beamer}
\usetheme{CambridgeUS}
\usepackage{xeCJK}
\usepackage{amsmath, amssymb, bm, mathtools}
\usepackage{color}

% 配色（蒋智超老师风格）
\definecolor{gold}{RGB}{255,192,0}    % 金黄强调 #FFC000
\definecolor{deepblue}{RGB}{5,99,193}  % 深蓝标题 #0563C1
\definecolor{grayblue}{RGB}{68,84,106} % 灰蓝正文 #44546A

\usecolortheme{default}
\setbeamercolor{title}{fg=deepblue}
\setbeamercolor{subtitle}{fg=grayblue}
\setbeamercolor{item}{fg=gold}

\begin{document}

% ============================================================
% Slide 1: 封面
% ============================================================
\begin{frame}
  \titlepage
\end{frame}

% ============================================================
% Slide 2: Motivation - 为什么需要 Gauss-Markov？
% ============================================================
\begin{frame}{为什么需要 Gauss-Markov？}
  \begin{block}{核心问题}
    OLS 估计凭什么``最好''？
  \end{block}
  \vspace{1em}
  \begin{itemize}
    \item 有没有比 OLS 更好的估计？
    \item 第 3 章 OLS 是纯代数运算，没有随机假设谈不上``最优''
  \end{itemize}
  \vspace{1em}
  \textbf{直觉说明：} Without any stochastic assumptions, the OLS in Chapter 3 is purely algebraic. If we want to discuss the statistical properties of OLS, we must invoke some statistical modeling assumptions.
\end{frame}

% ============================================================
% Slide 3: Motivation - 从代数到统计
% ============================================================
\begin{frame}{从代数到统计：引入随机假设的必要性}
  \begin{block}{核心观点}
    没有随机假设谈不上``最优''
  \end{block}
  \vspace{1em}
  \begin{itemize}
    \item 第 3 章 OLS 是纯代数运算
    \item Gauss-Markov 模型给出了讨论统计性质的基础
    \item 需要引入误差项 $\varepsilon$ 的前二阶矩假设
  \end{itemize}
  \vspace{1em}
  \textbf{直觉说明：} Gauss-Markov 模型 Assumption \S 4.1 给出随机假设，使得讨论估计量的均值、方差、协方差成为可能。
\end{frame}

% ============================================================
% Slide 4: 章节路线图
% ============================================================
\begin{frame}{章节路线图}
  \begin{block}{本章知识结构}
    \begin{enumerate}
      \item Gauss-Markov 模型假设
      \item OLS 估计量的性质（均值、方差）
      \item Gauss-Markov 定理（核心）
      \item 定理的直观解释与证明思路
    \end{enumerate}
  \end{block}
  \vspace{1em}
  \textbf{直觉说明：} 本章从模型假设出发，逐步建立 OLS 的统计性质，最终导出 Gauss-Markov 定理。
\end{frame}

% ============================================================
% Slide 5: Gauss-Markov 模型假设 (Assumption 4.1)
% ============================================================
\begin{frame}{Gauss-Markov 模型假设}
  \begin{block}{Assumption 4.1}
    \begin{align}
      Y &= X\beta + \varepsilon \quad \text{(线性形式)} \\
      E(\varepsilon) &= 0 \quad \text{(误差均值零)} \\
      \cov(\varepsilon) &= \sigma^2 I_n \quad \text{(同方差、不相关)} \\
      X &\text{ 固定且列线性无关}
    \end{align}
  \end{block}
  \vspace{1em}
  未知参数为 $(\beta, \sigma^2)$。

  \vspace{1em}
  \textbf{直觉说明：} $X$ 固定非本质（可条件于 $X$），关键是误差的前二阶矩。
\end{frame}

% ============================================================
% Slide 6: 个体水平视角
% ============================================================
\begin{frame}{个体水平视角：$y_i = x_i^\top \beta + \varepsilon_i$}
  \begin{block}{逐观测形式}
    \begin{align}
      E(\varepsilon_i) &= 0 \\
      \var(\varepsilon_i) &= \sigma^2 \quad \text{(同方差)} \\
      \cov(\varepsilon_i, \varepsilon_j) &= 0, \quad i \neq j \quad \text{(不相关)}
    \end{align}
  \end{block}
  \vspace{1em}
  \textbf{直觉说明：} 向量形式展开为 $n$ 个方程，每个误差独立满足相同的前二阶矩。
\end{frame}

% ============================================================
% Slide 7: 同方差假设的含义与重要性
% ============================================================
\begin{frame}{同方差假设的含义与重要性}
  \begin{block}{homoskedasticity}
    homoskedasticity = 相同方差（词源：homo = 相同，skedasis = 散布）
  \end{block}
  \vspace{1em}
  \begin{itemize}
    \item 误差项 $\varepsilon_i$ 具有相同方差 $\sigma^2$
    \item 异方差情形需要加权最小二乘（见第 19 章）
    \item 同方差假设是 Gauss-Markov 定理成立的关键
  \end{itemize}
  \vspace{1em}
  \textbf{直觉说明：} The assumption that the error terms have the same variance $\sigma^2$ is called homoskedasticity. The critiques on the assumptions aside, I will derive the properties of $\hat{\beta}$ under the Gauss-Markov model.
\end{frame}

% ============================================================
% Slide 8: OLS 估计量：矩阵形式
% ============================================================
\begin{frame}{OLS 估计量：矩阵形式}
  \begin{block}{OLS 估计量}
    \begin{align}
      \hat{\beta} = (X^\top X)^{-1} X^\top Y
    \end{align}
  \end{block}
  \vspace{1em}
  \begin{itemize}
    \item $\hat{\beta}$ 是 $Y$ 的线性函数
    \item 仅依赖矩阵运算
  \end{itemize}
  \vspace{1em}
  \textbf{直觉说明：} 令 $A = (X^\top X)^{-1} X^\top$，则 $\hat{\beta} = AY$，$A$ 不依赖 $Y$，所以 OLS 是线性估计量。
\end{frame}

% ============================================================
% Slide 9: Theorem 4.1 - OLS 估计量的均值与方差
% ============================================================
\begin{frame}{OLS 估计量的均值与方差}
  \begin{block}{Theorem 4.1 (无偏性)}
    在 Gauss-Markov 模型 Assumption 4.1 下：
    \begin{align}
      E(\hat{\beta}) = \beta
    \end{align}
  \end{block}
  \vspace{1em}
  \begin{block}{协方差矩阵}
    \begin{align}
      \cov(\hat{\beta}) = \sigma^2 (X^\top X)^{-1}
    \end{align}
  \end{block}
  \vspace{1em}
  \textbf{直觉说明：} OLS 估计量 $\hat{\beta}$ 是 $\beta$ 的无偏估计，其协方差矩阵完全由 $\sigma^2$ 和 $X$ 决定。
\end{frame}

% ============================================================
% Slide 10: Gauss-Markov 定理（核心）
% ============================================================
\begin{frame}{Gauss-Markov 定理（核心）}
  \begin{block}{Gauss-Markov Theorem 4.2}
    在 Gauss-Markov 模型 Assumption 4.1 下，$\hat{\beta}$ 是 BLUE（Best Linear Unbiased Estimator）。
  \end{block}
  \vspace{1em}
  \begin{block}{条件 C1 与 C2}
    对任意线性估计量 $\tilde{\beta} = AY$（$A$ 不依赖 $Y$）满足 $E(\tilde{\beta}) = \beta$（对所有 $\beta$ 无偏），有：
    \begin{align}
      \cov(\tilde{\beta}) \succeq \cov(\hat{\beta})
    \end{align}
  \end{block}
  \vspace{1em}
  \textbf{直觉说明：} We write $M_1 \succeq M_2$ if $M_1 - M_2$ is positive semi-definite.
\end{frame}

% ============================================================
% Slide 11: BLUE 的含义：协方差矩阵序
% ============================================================
\begin{frame}{BLUE 的含义：协方差矩阵序}
  \begin{block}{协方差序的等价定义}
    $\cov(\tilde{\beta}) \succeq \cov(\hat{\beta})$ 等价于：对任意 $c \in \mathbb{R}^p$，
    \begin{align}
      \var(c^\top \tilde{\beta}) \geq \var(c^\top \hat{\beta})
    \end{align}
  \end{block}
  \vspace{1em}
  \begin{itemize}
    \item 每个坐标分量 $\hat{\beta}_j$ 方差最小
    \item 对任意线性组合 $c^\top \tilde{\beta}$ 方差不小于 $c^\top \hat{\beta}$
  \end{itemize}
  \vspace{1em}
  \textbf{直觉说明：} So the OLS estimator has a smaller variance than other estimators for all coordinates.
\end{frame}

% ============================================================
% Slide 12: 为什么只比较线性估计量？
% ============================================================
\begin{frame}{为什么只比较线性估计量？}
  \begin{block}{线性约束的合理性}
    \begin{itemize}
      \item $A$ 可以是 $X$ 的任意非线性函数，线性约束已包含极广的估计量类
      \item 无偏性是自然要求
      \item 在许多现代应用中，有偏估计（如 Ridge、Lasso）方差更小
    \end{itemize}
  \end{block}
  \vspace{1em}
  \textbf{直觉说明：} Why do we restrict the estimator to be linear? The class of linear estimator is actually quite large because $A$ can be any nonlinear function of $X$.
\end{frame}

% ============================================================
% Slide 13: OLS 投影几何直观
% ============================================================
\begin{frame}{OLS 投影几何直观}
  \begin{block}{投影矩阵与分解}
    \begin{align}
      H &= X(X^\top X)^{-1} X^\top \quad \text{(投影矩阵)} \\
      \hat{Y} &= HY \quad \text{是 $Y$ 到 $X$ 列空间的投影} \\
      \hat{\varepsilon} &= (I_n - H)Y \quad \text{与列空间正交}
    \end{align}
  \end{block}
  \vspace{1em}
  OLS 将 $Y$ 分解为拟合值 $+$ 残差：$Y = \hat{Y} + \hat{\varepsilon}$

  \vspace{1em}
  \textbf{直觉说明：} $H$ 和 $I_n - H$ 是投影矩阵，$HX = X$，$(I_n - H)X = 0$，两者正交：$H(I_n - H) = 0$。
\end{frame}

% ============================================================
% Slide 14: 拟合值与残差的分布性质 (Theorem 4.2)
% ============================================================
\begin{frame}{拟合值与残差的分布性质}
  \begin{block}{Theorem 4.2}
    \begin{align}
      E(\hat{Y}) &= X\beta, \quad E(\hat{\varepsilon}) = 0 \\
      \cov(\hat{Y}) &= \sigma^2 H \\
      \cov(\hat{\varepsilon}) &= \sigma^2(I_n - H)
    \end{align}
  \end{block}
  \vspace{1em}
  $\hat{Y}$ 与 $\hat{\varepsilon}$ 不相关。

  \vspace{1em}
  \textbf{直觉说明：} 注意 $\hat{y}_i$ 和 $\hat{\varepsilon}_i$ 自身不同观测间是相关的（协方差非零）。
\end{frame}

% ============================================================
% Slide 15: 正交 vs. 不相关：两个关键陈述的区别
% ============================================================
\begin{frame}{正交 vs. 不相关：两个关键陈述的区别}
  \begin{block}{陈述 (S1)}
    $\hat{Y}$ 与 $\hat{\varepsilon}$ 正交——代数事实，无需随机假设（OLS 投影性质）
  \end{block}
  \vspace{0.5em}
  \begin{block}{陈述 (S2)}
    $\hat{Y}$ 与 $\hat{\varepsilon}$ 不相关——随机陈述，需要 Gauss-Markov 假设
  \end{block}
  \vspace{1em}
  两者含义不同，不能混淆。

  \vspace{1em}
  \textbf{直觉说明：} S1 是纯代数结论（Lemma 4.1），S2 需要 $\cov(Y) = \sigma^2 I_n$ 才能推导。
\end{frame}

% ============================================================
% Slide 16: 误差椭圆可视化：``最佳''的含义
% ============================================================
\begin{frame}{误差椭圆可视化：``最佳''的含义}
  \begin{block}{几何直观}
    \begin{itemize}
      \item $\hat{\beta}$ 的置信椭圆
      \item 其他线性无偏估计的置信椭圆
      \item OLS 的置信椭圆最小（被其他所有椭圆包裹）
      \item $\Rightarrow$ $\hat{\beta}$ 在所有方向上方差最小
    \end{itemize}
  \end{block}
  \vspace{1em}
  \textbf{直觉说明：} 几何直观：协方差矩阵正定半定序意味着椭圆包含关系，OLS 的置信椭圆被所有其他线性无偏估计的置信椭圆包含。
\end{frame}

% ============================================================
% Slide 17: Gauss-Markov 定理证明思路
% ============================================================
\begin{frame}{Gauss-Markov 定理证明思路}
  \begin{block}{四步框架}
    \textbf{第一步：} OLS 满足无偏性条件 $\Rightarrow AX = I_p$（由 $E(AY) = AX\beta = \beta$ 对所有 $\beta$ 成立推得）

    \vspace{0.5em}
    \textbf{第二步：} 协方差分解
    \begin{align}
      \cov(\tilde{\beta}) = \cov(\hat{\beta}) + \cov(\tilde{\beta} - \hat{\beta})
    \end{align}

    \vspace{0.5em}
    \textbf{第三步：} 交叉协方差项 $\cov(\hat{\beta}, \tilde{\beta} - \hat{\beta}) = 0$

    \vspace{0.5em}
    \textbf{第四步：} $\cov(\tilde{\beta} - \hat{\beta}) \succeq 0 \Rightarrow \cov(\tilde{\beta}) \succeq \cov(\hat{\beta})$
  \end{block}
  \vspace{1em}
  \textbf{直觉说明：} 证明核心是利用无偏性条件 $AX = I_p$ 消去交叉协方差项中的待定矩阵 $A$，从而分解式只剩 $\cov(\tilde{\beta} - \hat{\beta}) \succeq 0$。
\end{frame}

% ============================================================
% Slide 18: 例子：简单线性回归的 BLUE
% ============================================================
\begin{frame}{例子：简单线性回归的 BLUE}
  \begin{block}{一元情形}
    \begin{align}
      y_i &= \alpha + \beta x_i + \varepsilon_i \\
      \var(\hat{\beta}) &= \frac{\sigma^2}{\sum (x_i - \bar{x})^2}
    \end{align}
  \end{block}
  \vspace{1em}
  任何线性无偏估计的 $\beta$ 系数方差都不小于此值。

  \vspace{1em}
  \textbf{直觉说明：} 设计越分散（$\sum (x_i - \bar{x})^2$ 越大），OLS 估计越精确。
\end{frame}

% ============================================================
% Slide 19: 例子：均值的 BLUE
% ============================================================
\begin{frame}{例子：均值的 BLUE}
  \begin{block}{均值估计问题}
    $y_i \sim$ 均值 $\mu$，方差 $\sigma^2$，互不相关
  \end{block}
  \vspace{1em}
  线性估计 $\hat{\mu} = \sum a_i y_i$，无偏要求 $\sum a_i = 1$。

  \vspace{0.5em}
  \begin{block}{最优选择}
    $a_i = 1/n$（简单平均），方差 $\var(\hat{\mu}) = \sigma^2/n$ 为最小
  \end{block}
  \vspace{1em}
  \textbf{直觉说明：} 在无偏约束 $\sum a_i = 1$ 下，用 Lagrange 乘子法最小化 $\var(\hat{\mu}) = \sigma^2 \sum a_i^2$，最优解为均匀权重。
\end{frame}

% ============================================================
% Slide 20: 例子：Gauss-Markov 预测定理
% ============================================================
\begin{frame}{例子：Gauss-Markov 预测定理}
  \begin{block}{Theorem (Gauss-Markov for Prediction)}
    $\hat{Y} = X\hat{\beta}$ 是 $X\beta$ 的最佳线性无偏预测
  \end{block}
  \vspace{1em}
  适用于任何线性预测 $\tilde{Y} = \tilde{H} Y$。

  \vspace{0.5em}
  \begin{block}{核心不等式}
    \begin{align}
      \cov(\tilde{Y}) \succeq \cov(\hat{Y})
    \end{align}
  \end{block}
  条件：无偏 $E(\tilde{Y}) = X\beta$，线性 $\tilde{Y} = \tilde{H} Y$。

  \vspace{1em}
  \textbf{直觉说明：} 定理可推广到预测情形：OLS 预测 $\hat{Y}$ 在协方差矩阵序意义下优于所有线性无偏预测量。
\end{frame}

% ============================================================
% Slide 21: 总结
% ============================================================
\begin{frame}{总结}
  \begin{block}{Gauss-Markov 定理核心要点}
    \begin{enumerate}
      \item \textbf{Gauss-Markov 模型}：线性 $+$ 同方差 $+$ 不相关
      \item \textbf{OLS $\hat{\beta}$ 是 BLUE}：$\cov(\tilde{\beta}) \succeq \cov(\hat{\beta})$
      \item \textbf{核心不等式}：$\cov(\tilde{\beta}) - \cov(\hat{\beta}) = \cov(\tilde{\beta} - \hat{\beta}) \succeq 0$
    \end{enumerate}
  \end{block}
  \vspace{1em}
  \begin{block}{三个局限性}
    \begin{itemize}
      \item 不谈非线性估计（Ridge、Lasso 等有偏估计在现代高维问题中方差更小）
      \item 不谈正态假设（Normal 模型下 OLS 还具有 MLE 性质）
      \item 不谈稳健性（违背假设时的表现）
    \end{itemize}
  \end{block}
  \vspace{1em}
  \textbf{直觉说明：} Gauss-Markov 定理是线性估计理论的基础，核心洞察是：在唯一指定``线性 $+$ 无偏''后，OLS 自动最优，无需进一步选择。
\end{frame}

% ============================================================
% Slide 22: 参考文献
% ============================================================
\begin{frame}{参考文献}
  \begin{thebibliography}{99}
    \bibitem{PengDing2024} Peng Ding. \textit{A First Course in Causal Inference}. 2024.
    \bibitem{Sheather2009} Sheather, S.J. \textit{A Modern Approach to Regression with R}. Springer, 2009.
    \bibitem{McCulloch1985} McCulloch, J.H. \textit{On Heteroscedasticity}. Econometrica, 1985.
  \end{thebibliography}
\end{frame}

\end{document}
